\section{总结与延伸}

\subsection{讲者的核心总结}

Justin Johnson 在课程结尾强调了以下关键信息：

\begin{enumerate}
    \item \textbf{GPU 计算能力 1000 倍增长}是过去十年 AI 能力爆发的核心驱动力
    \item \textbf{内存层次结构}（从 SRAM 到跨 Pod 网络）决定了算法设计
    \item \textbf{数据并行 + FSDP + 激活检查点}可以覆盖大多数实际场景（$\leq$1,000 GPU）
    \item \textbf{MFU 是优化分布式训练的北极星指标}
    \item 真正的大规模训练（$>$10,000 GPU）需要 4D 并行 + 拓扑感知 + 自动容错
\end{enumerate}

\subsection{全课知识图谱}

\begin{center}
\begin{tikzpicture}[node distance=0.8cm, every node/.style={draw, rounded corners, minimum width=3.5cm, minimum height=0.7cm, font=\small}]
\node (hw) {GPU 硬件架构};
\node (cluster) [below=of hw] {GPU 集群层次结构};
\node (dp) [below=of cluster] {DP / FSDP / HSDP};
\node (adv) [below=of dp] {CP / PP / TP};
\node (nd) [below=of adv] {4D 并行 + MFU 优化};
\draw[->, thick] (hw) -- (cluster);
\draw[->, thick] (cluster) -- (dp);
\draw[->, thick] (dp) -- (adv);
\draw[->, thick] (adv) -- (nd);
\node[draw=none, right=1cm of hw, text width=5cm, font=\scriptsize] {SM、Tensor Core、混合精度};
\node[draw=none, right=1cm of cluster, text width=5cm, font=\scriptsize] {Server/Rack/Pod/Cluster、带宽层次};
\node[draw=none, right=1cm of dp, text width=5cm, font=\scriptsize] {梯度 All-Reduce、权重分片、激活检查点};
\node[draw=none, right=1cm of adv, text width=5cm, font=\scriptsize] {Ring/Ulysses Attention、Micro-batch、Block MatMul};
\node[draw=none, right=1cm of nd, text width=5cm, font=\scriptsize] {拓扑映射、容错、Llama 3 案例};
\end{tikzpicture}
\end{center}

\subsection{关键 Takeaways}

\begin{importantbox}{五条核心原则}
\begin{enumerate}
    \item \textbf{Tensor Core 是 GPU 的灵魂}：确保你的计算尽量走 Tensor Core（使用 16-bit 精度、矩阵乘法为主的架构）
    \item \textbf{带宽随距离急剧衰减}：算法设计必须尊重 GPU 集群的层次化带宽结构
    \item \textbf{DP/FSDP/HSDP 覆盖大多数场景}：这是你最可能在实践中使用的技术
    \item \textbf{MFU 是唯一的优化指标}：面对众多并行策略和超参数时，MFU 告诉你方向
    \item \textbf{容错是大规模训练的必需品}：在数万 GPU 上训练数月，硬件故障不是``如果''的问题，而是``何时''的问题
\end{enumerate}
\end{importantbox}

\subsection{拓展阅读}

\begin{itemize}
    \item Llama 3 技术报告：\url{https://arxiv.org/abs/2407.21783}
    \item PyTorch FSDP 文档：\url{https://pytorch.org/docs/stable/fsdp.html}
    \item Megatron-LM 论文（张量并行）：Shoeybi et al., \url{https://arxiv.org/abs/1909.08053}
    \item Ring Attention：Li et al., \url{https://arxiv.org/abs/2310.01889}
    \item GPipe（流水线并行）：Huang et al., \url{https://arxiv.org/abs/1811.06965}
    \item ZeRO（FSDP 的前身）：Rajbhandari et al., \url{https://arxiv.org/abs/1910.02054}
\end{itemize}

\end{document}
